\chapter{Testing di applicazioni web}

Abbiamo imparato a progettare e realizzare applicazioni web moderne e sicure.
Manca un ultimo tassello prima di poter consegnare applicazioni di
\textbf{qualità}: il testing. Testare è fondamentale per intercettare quanti più
difetti possibile \textbf{prima} che raggiungano la produzione, dove possono
costare molto cari.

\section{Richiami}

\dfn{Verifica del software}{
  L'insieme delle pratiche volte ad accertare che un software soddisfi una
  specifica. Si distingue in \term{verifica dinamica}, che comporta
  l'esecuzione del programma (il \textbf{testing}), e \term{verifica statica},
  che non la comporta (revisione del codice, analisi statica, verifica
  automatica).}

Un \term{test} è una sequenza di azioni progettate per valutare un particolare
aspetto o funzionalità di un software; il software sotto test si indica
comunemente con \term{SUT} (\textit{Software Under Test}). Un test consiste
generalmente in tre fasi:

\begin{itemize}
  \item \term{Arrange}: assicurarsi che le precondizioni richieste siano
        soddisfatte;
  \item \term{Act}: eseguire una o più azioni;
  \item \term{Assert}: verificare che il SUT si sia comportato come previsto.
\end{itemize}

\subsection{I livelli di testing}

\begin{center}
\renewcommand{\arraystretch}{1.4}
\begin{tabular}{@{}l@{\hspace{1.2cm}}>{\raggedright\arraybackslash}p{8.6cm}@{}}
  \textbf{Livello} & \textbf{Bersaglio} \\
  \term{Unit testing}        & una singola unità (una classe, un metodo) \\
  \term{Integration testing} & la collaborazione fra unità diverse \\
  \term{System testing}      & il sistema nel suo complesso, rispetto alla
                               specifica \\
  \term{End-to-End testing}  & il sistema dal punto di vista dei suoi utenti
                               finali \\
\end{tabular}
\end{center}

\begin{figure}[H]
  \centering
  \begin{tikzpicture}[font=\Lato\scriptsize,
      lvl/.style={line width=0.8pt, align=center}]

    \fill[red-gradient-1!50!white, draw=red-gradient-6, line width=0.8pt]
      (-1.0,3.2) -- (1.0,3.2) -- (0.7,2.4) -- (-0.7,2.4) -- cycle;
    \node at (0,2.85) {Manuale};

    \fill[orange-gradient-1!50!white, draw=orange-gradient-6, line width=0.8pt]
      (-1.6,2.35) -- (1.6,2.35) -- (1.3,1.55) -- (-1.3,1.55) -- cycle;
    \node at (0,1.98) {E2E};

    \fill[yellow-gradient-2!50!white, draw=yellow-gradient-6, line width=0.8pt]
      (-2.4,1.5) -- (2.4,1.5) -- (2.1,0.7) -- (-2.1,0.7) -- cycle;
    \node at (0,1.13) {Test di sistema};

    \fill[azure-gradient-1!60!white, draw=azure-gradient-6, line width=0.8pt]
      (-3.3,0.65) -- (3.3,0.65) -- (3.0,-0.15) -- (-3.0,-0.15) -- cycle;
    \node at (0,0.28) {Test di integrazione};

    \fill[green-gradient-1!60!white, draw=green-gradient-6, line width=0.8pt]
      (-4.3,-0.2) -- (4.3,-0.2) -- (4.0,-1.1) -- (-4.0,-1.1) -- cycle;
    \node at (0,-0.62) {Test di unità};

    \draw[-{Stealth[length=6pt]}, line width=0.9pt, draw=neutral!45!black]
      (-5.2,-1.1) -- (-5.2,3.2);
    \node[font=\Lato\scriptsize, text=neutral!35!black, rotate=90, anchor=south]
      at (-5.4,1.0) {più lenti e costosi, più vicini all'utente};

    \draw[-{Stealth[length=6pt]}, line width=0.9pt, draw=neutral!45!black]
      (5.2,3.2) -- (5.2,-1.1);
    \node[font=\Lato\scriptsize, text=neutral!35!black, rotate=-90, anchor=south]
      at (5.4,1.0) {più numerosi, più veloci ed economici};
  \end{tikzpicture}
  \caption{La piramide del testing.}
  \label{fig:testing-pyramid}
\end{figure}

\subsection{Le difficoltà specifiche del Web}

Le applicazioni web, come ogni software, vanno testate. Emergono però alcune
difficoltà peculiari:

\begin{itemize}
  \item nel \textbf{unit testing} bisogna fare i conti con il codice
        \textbf{asincrono} e con le dipendenze da \textbf{servizi esterni};
  \item nel \textbf{testing end-to-end} bisogna fare i conti con
        \textbf{fragilità} e \textbf{flakiness}.
\end{itemize}

\section{Unit testing con Angular}

Angular fornisce tutto il necessario per testare le applicazioni usando
\term{Jasmine}. Per impostazione predefinita il codice di test si scrive in file
\term{spec} (specification), e la CLI ne genera degli scheletri: si collocano
nella stessa directory del componente o servizio che testano, con la convenzione
di aggiungere \code{.spec} al nome.

\begin{lstlisting}[style=ts, name=calculator.service.spec.ts]
import { TestBed } from '@angular/core/testing';
import { CalculatorService } from './calculator.service';

describe('CalculatorService', () => {
  let service: CalculatorService;

  beforeEach(() => {
    TestBed.configureTestingModule({});
    service = TestBed.inject(CalculatorService);
  });

  it('should be created', () => {
    expect(service).toBeTruthy();
  });

  it('should sum two positive numbers correctly', () => {
    expect(service.sum(42, 58)).toEqual(100);
  });
});
\end{lstlisting}

\begin{center}
\renewcommand{\arraystretch}{1.4}
\begin{tabular}{@{}l@{\hspace{1.2cm}}>{\raggedright\arraybackslash}p{8.6cm}@{}}
  \textbf{Elemento} & \textbf{Ruolo} \\
  \code{describe()} & crea una \textbf{suite} di specifiche, raggruppando test
                      correlati. È l'equivalente di una classe di test JUnit. \\
  \code{it()}       & definisce una singola specifica: l'equivalente di un
                      metodo annotato \code{@Test}. \\
  \code{TestBed}    & la principale utilità di testing di Angular: fornisce i
                      metodi per creare componenti e servizi nei test di
                      unità. \\
\end{tabular}
\end{center}

Il comando \code{ng test} compila l'applicazione e lancia il test runner
\term{Karma}, che raccoglie tutti i file spec, li esegue e produce un rapporto.

\section{Dipendenze e test double}

Tipicamente i servizi e i componenti che testiamo dipendono da altri servizi o
componenti: il servizio \code{RestBackend} dipende da \code{HttpClient}, il
componente \code{TodoPage} dipende da \code{RestBackend}.

In questi scenari il unit testing diventa più delicato. Se il componente
\code{TodoPage} non mostra gli elementi attesi, dipende da un bug nel
componente? Nel servizio? Nel backend REST esterno? O da una combinazione?

Il problema va oltre la localizzazione del difetto. In molti casi non
\textbf{vogliamo} usare il codice di produzione reale durante i test. Si pensi
a un \code{TemperatureService} che restituisce la temperatura attuale a Napoli
interrogando una API esterna: come si testa che cosa accade quando la
temperatura scende sotto zero? Aspettando qualche anno una giornata
eccezionalmente fredda?

\dfn{Test double}{
  Un buon test di unità dovrebbe testare l'unità \textbf{in isolamento}: il suo
  esito non dovrebbe dipendere da altre unità. I \term{test double} sono
  sostituti degli oggetti di produzione, usati nei test al posto di questi.}

Invece del \code{TemperatureService} reale si usa un doppio che restituisce
sempre $-42$\,°C.

\subsection{Iniezione delle dipendenze e double}

Il meccanismo di iniezione delle dipendenze è un supporto prezioso: durante i
test si può configurare l'iniettore perché risolva le dipendenze usando i test
double invece degli oggetti reali. È il motivo per cui, nel capitolo su Angular,
insistevamo sul fatto che un componente non dovrebbe mai costruire da sé i
propri servizi.

Scrivere i double a mano richiede però molto codice ripetitivo. I framework di
testing vengono in aiuto: Jasmine supporta il concetto di \term{spy}, un test
double che può simulare qualunque funzione.

\begin{lstlisting}[style=ts, name=spy.ts]
let todos = [{ id: 1, todo: "foo" }, { id: 2, todo: "bar" }];

restBackendSpy = jasmine.createSpyObj('RestBackendService', ['getTodos']);
restBackendSpy.getTodos.and.returnValue(of(todos));
\end{lstlisting}

Il codice crea uno spy del \code{RestBackendService} simulandone il metodo
\code{getTodos()}, e lo configura perché restituisca sempre un Observable con
una data lista di due elementi. Usando lo spy al posto del servizio reale, il
test non dipende né dal servizio né dalla API REST esterna.

\begin{esempio}{Testare un componente}
\begin{lstlisting}[style=ts, name=todo-page.component.spec.ts]
it('should properly fetch to-do items', () => {
  let todos = [{ id: 1, todo: "foo" }, { id: 2, todo: "bar" }];

  restBackendSpy = jasmine.createSpyObj('RestBackendService', ['getTodos']);
  restBackendSpy.getTodos.and.returnValue(of(todos));
  component.restService = restBackendSpy;

  component.fetchTodos();
  fixture.detectChanges();

  expect(component.todos.length).toEqual(2);

  const element: DebugElement = fixture.debugElement;
  const list = element.query(By.css('ul'));
  const content = list.nativeElement.textContent;

  expect(content).toContain('foo');
  expect(content).toContain('bar');
});
\end{lstlisting}

Si noti \code{fixture.detectChanges()}: forza il ciclo di change detection
descritto nel capitolo 20, senza il quale il template non sarebbe ancora
aggiornato al momento delle asserzioni.
\end{esempio}

\subsection{Copertura del codice}

\begin{lstlisting}[style=shell]
ng test --no-watch --code-coverage
\end{lstlisting}

\begin{console}[Terminale]
Chrome 121.0.0.0 (Windows 10): Executed 18 of 18 SUCCESS
TOTAL: 18 SUCCESS

=========================== Coverage summary ===========================
Statements   : 45.45% ( 75/165 )
Branches     : 21.21% ( 7/33 )
Functions    : 43.39% ( 23/53 )
Lines        : 42.58% ( 66/155 )
========================================================================
\end{console}

Viene generato anche un rapporto HTML nella directory \code{/coverage}, dove si
può ispezionare ciascun file e vedere quali righe sono state coperte dai test.

\warning{La copertura non è la qualità}{
  La copertura è un \textbf{possibile indicatore} di quanto bene si è testata
  un'applicazione, non una misura della qualità dei test. Il 100\% di copertura
  si può ottenere anche con test che non verificano nulla: eseguono il codice
  senza mai controllarne il risultato. È utile soprattutto al contrario, per
  individuare le parti \textbf{mai} eseguite.}

\section{Testing end-to-end}

\dfn{Testing end-to-end}{
  Il \term{testing end-to-end} (E2E) mira a testare il sistema nel suo
  complesso, dal punto di vista dei suoi utenti finali. Ogni test replica un
  flusso d'uso realistico, interagendo con le interfacce \textbf{esterne} del
  sistema, come farebbero gli utenti.}

\subsection{E2E per una API REST}

Quando l'applicazione è una API REST, gli \textbf{endpoint} sono l'interfaccia
esterna e gli utenti finali sono \textbf{altri programmi} che inviano richieste
HTTP. I test E2E devono quindi inviare richieste HTTP e verificare che le
risposte siano corrette rispetto alla specifica.

\begin{center}
\renewcommand{\arraystretch}{1.4}
\begin{tabular}{@{}c@{\hspace{0.6cm}}>{\raggedright\arraybackslash}p{10.4cm}@{}}
  \multicolumn{2}{@{}l}{\color{primary!80!black}\textbf{Scenario: login riuscito}} \\
  1 & inviare una \code{POST} a \code{/auth} con corpo
      \code{\{usr: "gerry", pwd: "sussman"\}} \\
  2 & verificare che la risposta abbia stato \code{200 OK} \\
  3 & verificare che il valore restituito sia un JSON di tipo
      \code{\{token: string\}} \\
  4 & verificare che \code{response.token} sia un token JWT valido \\
\end{tabular}
\end{center}

\begin{center}
\renewcommand{\arraystretch}{1.4}
\begin{tabular}{@{}c@{\hspace{0.6cm}}>{\raggedright\arraybackslash}p{10.4cm}@{}}
  \multicolumn{2}{@{}l}{\color{primary!80!black}\textbf{Scenario: inserimento di
    un nuovo elemento}} \\
  1 & inviare una \code{POST} a \code{/auth} con le credenziali \\
  2 & verificare lo stato \code{200 OK}, estrarre e salvare il token \\
  3 & inviare una \code{POST} a \code{/todos} con corpo
      \code{\{todo: "foobar", done: false\}} e header
      \code{Authorization: Bearer <TOKEN>} \\
  4 & verificare che la risposta abbia stato \code{200 OK} \\
  5 & verificare che il corpo contenga
      \code{\{todo: "foobar", done: false\}} \\
\end{tabular}
\end{center}

\subsection{E2E per un'applicazione web}

Quando l'applicazione è tradizionale o è una SPA, le \textbf{pagine} sono
l'interfaccia esterna e gli utenti finali sono \textbf{persone} che usano un
browser. I test devono quindi interagire con le pagine e verificare che
cambino correttamente in conseguenza delle interazioni.

\begin{center}
\renewcommand{\arraystretch}{1.4}
\begin{tabular}{@{}c@{\hspace{0.6cm}}>{\raggedright\arraybackslash}p{10.4cm}@{}}
  \multicolumn{2}{@{}l}{\color{primary!80!black}\textbf{Scenario: login riuscito}} \\
  1 & visitare la pagina di login \\
  2 & inserire \guillemotleft gerry\guillemotright\ nel campo del nome utente \\
  3 & inserire \guillemotleft sussman\guillemotright\ nel campo della password \\
  4 & fare clic sul pulsante \guillemotleft Login\guillemotright \\
  5 & verificare che l'utente venga reindirizzato alla homepage \\
  6 & verificare che la barra di navigazione contenga \guillemotleft Hi, Gerry\guillemotright \\
\end{tabular}
\end{center}

Che sia manuale o automatizzato, il testing E2E segue sempre lo stesso schema:

\begin{lstlisting}[style=plain, name=pseudocodice]
esegui(suite di test):
  per ogni scenario nella suite:
    per ogni passo dello scenario:
      seleziona l'elemento con cui interagire (pulsante, campo, ...)
      interagisci con esso (clic, riempi, verifica una proprieta', ...)
\end{lstlisting}

\subsection{Manuale o automatizzato}

\begin{center}
\renewcommand{\arraystretch}{1.4}
\begin{tabular}{@{}>{\raggedright\arraybackslash}p{6.3cm}@{\hspace{1.0cm}}>{\raggedright\arraybackslash}p{6.3cm}@{}}
  \textbf{Manuale} & \textbf{Automatizzato} \\
  A persone viene dato un copione con l'elenco dettagliato dei passi; lo
  seguono e replicano segnalando ogni problema.
  & Si sviluppa software di test che simula automaticamente gli scenari
    d'uso. \\
  Non richiede competenze di programmazione.
  & Richiede competenze di programmazione e di testing. \\
  Attività lenta, noiosa e soggetta a errori. Non scala: va ripetuta a ogni
  nuova versione, su tutti i browser principali, su dispositivi diversi.
  & Costi iniziali più alti, ma le suite si possono rieseguire molte volte con
    sforzo aggiuntivo minimo. \\
  & Il codice di test va \textbf{mantenuto} man mano che l'applicazione evolve, e
    i test possono essere fragili o \textit{flaky}. \\
\end{tabular}
\end{center}

I test E2E automatizzati sono in sostanza codice che usa librerie dedicate per
\textbf{controllare da remoto un browser}: navigare verso URL, individuare
elementi e interagirci. Queste librerie sono utili anche oltre il testing, per
esempio per lo \textit{scraping} e il \textit{crawling}.

\section{Individuare gli elementi}

Selezionare gli elementi con cui interagire è una parte cruciale del testing E2E.
Esistono tre \term{strategie di localizzazione}.

\begin{center}
\renewcommand{\arraystretch}{1.45}
\begin{tabular}{@{}l@{\hspace{0.9cm}}>{\raggedright\arraybackslash}p{8.8cm}@{}}
  \textbf{Strategia} & \textbf{Come funziona} \\
  \textbf{Coordinate assolute}
    & identifica gli elementi in base alla posizione sullo schermo:
      \code{click(200, 300)} \\
  \textbf{Localizzatori visuali}
    & identifica gli elementi in base al loro \textbf{aspetto}, con algoritmi di
      confronto fra immagini \\
  \textbf{Localizzatori di layout}
    & identifica gli elementi in base a proprietà strutturali, per esempio con i
      \textbf{selettori CSS}: \code{click(".item button.buy")} \\
\end{tabular}
\end{center}

Torna qui utile tutto quello che abbiamo studiato sui selettori CSS nel capitolo
3: sono di gran lunga la strategia più usata nella pratica.

\section{Fragilità e flakiness}

\subsection{Fragilità}

\dfn{Fragilità}{
  I test E2E sono spesso \term{fragili} rispetto all'evoluzione
  dell'applicazione: anche modifiche minime possono rompere i localizzatori,
  rendendo i test incapaci di interagire con gli elementi corretti e richiedendo
  una manutenzione onerosa.}

\begin{center}
\renewcommand{\arraystretch}{1.45}
\begin{tabular}{@{}>{\raggedright\arraybackslash}p{5.0cm}@{\hspace{0.9cm}}>{\raggedright\arraybackslash}p{7.6cm}@{}}
  \textbf{Localizzatore} & \textbf{Perché si rompe} \\
  \code{click(100, 200)}
    & basta spostare il pulsante di qualche pixel e le coordinate non
      corrispondono più. \\
  Immagine del pulsante
    & basta cambiare il colore o il carattere e l'immagine non corrisponde
      più. \\
  \code{click(".form button:first-of-type")}
    & se si aggiunge un pulsante \guillemotleft Sign up\guillemotright\ \textbf{prima} di quello di
      login, il test farà clic su quello sbagliato --- e fallirà in modo
      silenzioso e fuorviante. \\
\end{tabular}
\end{center}

\info{Come si scrivono buoni localizzatori}{
  Un buon localizzatore non deve essere né troppo generico né troppo specifico:
  se è troppo generico rischia di selezionare un elemento diverso da quello
  voluto, se è troppo specifico anche modifiche minime lo rompono.

  I localizzatori a coordinate sono i più fragili in assoluto. Quelli visuali
  sono robusti rispetto ai cambiamenti di disposizione ma fragili rispetto
  all'aspetto; quelli di layout sono robusti rispetto ai cambiamenti puramente
  visivi ma più fragili rispetto ai cambiamenti di struttura. Nella pratica
  l'esperienza conta molto.}

\subsection{Flakiness}

\dfn{Flakiness}{
  I test E2E tendono a essere \term{flaky}, cioè a comportarsi in modo
  incoerente: lo stesso test può passare o fallire in modo intermittente
  \textbf{a parità di condizioni}, senza alcuna modifica all'applicazione o al
  test.}

Questa imprevedibilità è spesso causata dal \textbf{non determinismo} legato ai
tempi di caricamento o al recupero dei dati: talvolta una risorsa esterna
impiega più tempo del solito, e il test prova ad accedere a elementi non ancora
disegnati, ottenendo un errore di \guillemotleft elemento non trovato\guillemotright.

La flakiness si mitiga con opportune \textbf{strategie di attesa}: non attese a
tempo fisso, ma attese su una condizione (\guillemotleft finché quell'elemento non è
visibile\guillemotright).

\subsection{Isolamento}

I test E2E dovrebbero essere eseguiti in \textbf{isolamento}:

\begin{itemize}
  \item nessuna propagazione dei fallimenti: un test che fallisce non deve
        influenzare quelli successivi;
  \item l'ordine di esecuzione non deve contare, così che i test possano girare
        in parallelo;
  \item ogni test deve poter essere eseguito da solo, il che rende molto più
        facile il debug.
\end{itemize}

Si ottiene in due modi: \textbf{ripartire da zero}, eseguendo ogni test in un
ambiente nuovo e pulito, oppure fare \textbf{pulizia fra un test e l'altro}.

\section{Playwright}

\begin{center}
\renewcommand{\arraystretch}{1.4}
\begin{tabular}{@{}l@{\hspace{1.0cm}}l@{\hspace{1.0cm}}>{\raggedright\arraybackslash}p{5.2cm}@{}}
  \textbf{Strumento} & \textbf{Dal} & \textbf{Note} \\
  \term{Selenium}   & 2009 & multi-browser \\
  \term{Cypress}    & 2014 & multi-browser \\
  \term{Puppeteer}  & 2017 & di Google, orientato a Chrome \\
  \term{Playwright} & 2020 & di Microsoft, multi-browser \\
\end{tabular}
\end{center}

\term{Playwright} è un framework open source di automazione e testing dei
browser, sviluppato da Microsoft. È multi-browser e multi-piattaforma, adotta un
approccio asincrono guidato dagli eventi, è scritto in TypeScript e ha API in
TypeScript, JavaScript, C\#, Java e Python.

\begin{lstlisting}[style=shell]
npm init playwright@latest
\end{lstlisting}

\subsection{Il primo test}

I test si definiscono con la funzione \code{test()}, che riceve un titolo e una
funzione asincrona contenente il codice.

\begin{lstlisting}[style=ts, name=example.spec.ts]
import { test, expect } from '@playwright/test';

test('has title', async ({ page }) => {
  await page.goto('https://playwright.dev/');
  await expect(page).toHaveTitle(/Playwright/);
});

test('get started link', async ({ page }) => {
  await page.goto('https://playwright.dev/');

  await page.getByRole('link', { name: 'Get started' }).click();

  await expect(page.getByRole('heading', { name: 'Installation' }))
    .toBeVisible();
});
\end{lstlisting}

\begin{console}[Terminale]
$ npx playwright test
Running 6 tests using 6 workers
  6 passed (7.1s)

$ npx playwright show-report
Serving HTML report at http://localhost:9323.
\end{console}

Per impostazione predefinita i test vengono eseguiti su \code{chromium} (motore
Blink), \code{firefox} (Gecko) e \code{webkit} (il motore di Safari e dei
dispositivi Apple). Playwright può anche emulare browser per tablet e telefoni,
e le piattaforme si configurano nel file \code{playwright.config.ts}.

\subsection{Isolamento e fixture}

Playwright ottiene l'isolamento \textbf{ripartendo da zero}: ogni test gira in un
ambiente pulito, non influenzato dai precedenti, con il proprio
\code{localStorage}, \code{sessionStorage} e i propri cookie. Questo si ottiene
usando un nuovo \code{BrowserContext} per ciascun test --- si può pensarlo come
un profilo in incognito.

L'ambiente di ciascun test è definito tramite \term{fixture}. L'argomento
\code{\{page\}} passato alla funzione di test dice a Playwright di preparare la
fixture \code{Page}, cioè di creare un nuovo contesto e una pagina al suo
interno --- l'astrazione di una scheda in un browser nuovo di zecca.

\subsection{Individuare, agire, verificare}

Tre parti sono cruciali nella scrittura di un test E2E automatizzato:
\textbf{individuare} gli elementi, \textbf{compiere} azioni e
\textbf{verificare} lo stato rispetto alle aspettative.

\begin{lstlisting}[style=ts, name=localizzatori.ts]
page.getByText();         // per contenuto testuale
page.getByLabel();        // un controllo di form, tramite la sua etichetta
page.getByPlaceholder();  // un input, tramite il suo placeholder
page.getByAltText();      // un elemento (di solito immagine), tramite alt
page.getByTitle();        // tramite l'attributo title
page.getByTestId();       // tramite l'attributo data-testid

// e' possibile usare anche selettori CSS e XPath
page.locator("#foo ul [data-foo]");
page.locator("//[@id='foo']//ul//*[@data-foo]");
\end{lstlisting}

\begin{lstlisting}[style=ts, name=azioni.ts]
page.locator("#input").clear();
page.locator("#input").fill("Foo");
page.locator("#input").press("Tab");
page.locator("#checkbox").check();
page.locator("#select-course").selectOption("web technologies");
page.locator("#button").click();
page.locator("#button").dblclick();
\end{lstlisting}

\begin{lstlisting}[style=ts, name=asserzioni.ts]
expect(page.locator("#foo")).toBeDefined();
expect(page.locator("#foo")).toBeDisabled();
expect(page.locator("#foo")).toBeEmpty();
expect(page.locator("#foo")).toContainText("foo");
expect(page.locator("#foo")).toHaveClass("bar");
expect(page.locator("#foo")).toBeVisible();
expect(page.locator("#foo")).toBeInViewport();
\end{lstlisting}

\info{Perché i localizzatori \guillemotleft semantici\guillemotright\ sono preferibili}{
  \code{getByRole}, \code{getByLabel} e \code{getByPlaceholder} individuano gli
  elementi per il loro \textbf{ruolo} nell'interfaccia, non per la loro
  posizione nell'albero. Sono quindi molto meno fragili di un selettore CSS
  legato alla struttura, e hanno un effetto collaterale utile: un test che non
  riesce a trovare un pulsante \guillemotleft per ruolo\guillemotright\ segnala spesso un problema
  reale di accessibilità.}

\section{Capture and Replay}

Scrivere i test a mano richiede tempo. Gli approcci \term{Capture and Replay}
(C\&R) permettono di generare automaticamente il codice dei test
\guillemotleft registrando\guillemotright\ le interazioni dell'utente: basta interagire con
l'applicazione per creare test rieseguibili.

Essendo generato automaticamente, il codice C\&R può essere subottimale rispetto
a fragilità e flakiness, ma resta un buon punto di partenza.

Playwright supporta il C\&R attraverso il \textit{Test Generator}. Il modo più
semplice è usare l'estensione Playwright per VS Code: nella scheda
\textit{Testing} si sceglie \guillemotleft Record new\guillemotright, si interagisce con il browser
che compare, e Playwright genera il codice necessario a replicare le azioni,
usando euristiche interne per scegliere i localizzatori migliori.

\begin{lstlisting}[style=ts, name=login.spec.ts (generato)]
test('should login with correct credentials', async ({ page }) => {
  await page.goto('http://localhost:4200/home');
  await page.getByRole('link', { name: 'Login' }).click();
  await page.getByPlaceholder('Your username').fill('luigi');
  await page.getByPlaceholder('********').fill('luigi');
  await page.getByRole('button', { name: 'Sign in' }).click();
  await expect(page.locator('#navbar-default')).toContainText('Welcome, luigi');
});
\end{lstlisting}

\section{Riferimenti}

\begin{itemize}
  \item \textbf{Angular Testing}. \\
        \urlx{https://angular.dev/guide/testing} \\
        Parti rilevanti: \textit{Overview}, \textit{Basics of testing
        components}, \textit{Component Testing Scenarios}, \textit{Testing
        services}, \textit{Code coverage}.
  \item \textbf{Playwright Docs}. \\
        \urlx{https://playwright.dev/docs/intro} \\
        Parti rilevanti: \textit{Getting started}, \textit{Guides > Best
        practices}.
\end{itemize}
